
\section{Positive transport and cyclic descent}

Let $t$ be the coordinate on a complex disk, oriented by increasing
argument. A positive meridian is the loop $t=\epsilon e^{2\pi iv}$,
$0\le v\le1$, with $\epsilon>0$. The integral monodromy of a family
along this loop is the map on the homology of the starting fibre
obtained by transporting cycles and identifying the ending fibre
with the starting fibre. Fix the cyclic change of coordinates $t=s^3$
and the generator $s\mapsto\zeta s$, where
$\zeta=e^{2\pi i/3}$. A deck transformation is an automorphism of
the pulled-back family covering this rotation and acting trivially
on its quotient.

\begin{lemma}\label{marked-local:lem:inverse-descent}
Suppose a smooth family over the $s$ disk has a deck transformation
$\rho$ of order three. Mark its integral first homology by parallel
continuation over the disk, and let $A$ be the matrix of $\rho$ in
that marking. The positive monodromy of the quotient over $t\ne0$
is $A^{-1}$.
\end{lemma}
\begin{proof}
The positive meridian lifts from $s=\epsilon^{1/3}$ to
$s=\zeta\epsilon^{1/3}$. A homology class remains constant in the
disk marking along this lifted path. To express its endpoint in the
starting fibre of the quotient one applies $\rho^{-1}$. This acts
by $A^{-1}$. A translation of a torus is homotopic to the identity
through translations, so the same argument applies to an affine
deck transformation on its first homology. It makes no assertion
that the affine translation itself vanishes.
\end{proof}

The inverse can be seen without a monodromy table. Consider the
projective elliptic family whose affine equation is
\begin{equation}\label{marked-local:eq:weierstrass}
 y^2=x^3-3t^3(t-1)x+2t^4(t-1)^2.
\end{equation}
Its origin is the point at infinity. For $s\ne0$, put
$t=s^3$, $x=s^4X$, and $y=s^6Y$. Dividing by $s^{12}$ gives
\begin{equation}\label{marked-local:eq:good-equation}
 Y^2=X^3-3s(s^3-1)X+2(s^3-1)^2.
\end{equation}
For an equation $Y^2=X^3+aX+b$ its discriminant is
$-16(4a^3+27b^2)$; it vanishes exactly when the cubic and its
derivative have a common root. Here it is
$1728(s^3-1)^3$, so the projective fibres are smooth on a
sufficiently small disk around zero. The central curve is
$Y^2=X^3+2$. The deck transformation that fixes $x,y,t$ is
\begin{equation}\label{marked-local:eq:actual-elliptic-deck}
 (X,Y,s)\longmapsto(\zeta^2X,Y,\zeta s).
\end{equation}
Substitution in \eqref{marked-local:eq:good-equation} verifies the transformation
on the whole smooth family, including the origin section.

On the central curve its pullback on the holomorphic differential
$dX/Y$ is multiplication by $\zeta^2$. Thus for any cycle $c$,
\[
 \int_{\rho_*c}\frac{dX}{Y}=\zeta^2\int_c\frac{dX}{Y}.
\]
In a period marking with basis $(1,\zeta)$ the central deck matrix
is consequently
$\left(\begin{smallmatrix}-1&1\\-1&0\end{smallmatrix}\right)$.
The positive return matrix is its inverse. This calculation fixes
both the disk orientation and the direction of action on periods.

For the four-torus, retain the following explicit marked lattice
in $\mathbb C^2$. Choose $\tau$ with $\operatorname{Im}\tau>0$ and set
\begin{equation}\label{marked-local:eq:period-lattice}
 e_1=(1,0),\quad e_2=(\zeta,0),\quad
 e_3=\tfrac13(2+\zeta,1),\quad e_4=(0,\tau),
 \qquad \Lambda=\sum_{j=1}^4\mathbb Ze_j.
\end{equation}
The vector $\delta=(0,1)=-2e_1-e_2+3e_3$ is an integral
period. The map multiplying the first coordinate by $\zeta$ and
fixing the second preserves $\Lambda$ and has matrix
\begin{equation}\label{marked-local:eq:positive-T}
 T=\begin{pmatrix}0&-1&-1&0\\1&-1&0&0\\0&0&1&0\\0&0&0&1\end{pmatrix}.
\end{equation}
Indeed its values on the four basis vectors are
$e_2,-e_1-e_2,e_3-e_1,e_4$. Its inverse is the integral map
\begin{equation}\label{marked-local:eq:inverse-T}
 A=T^{-1}=
 \begin{pmatrix}-1&1&-1&0\\-1&0&-1&0\\0&0&1&0\\0&0&0&1\end{pmatrix},
\end{equation}
with values $-e_1-e_2,e_1,e_3-e_1-e_2,e_4$.
In particular it fixes $\delta$ and $e_4$.

If \eqref{marked-local:eq:positive-T} is the prescribed positive monodromy in
this unchanged disk marking, Lemma~\ref{marked-local:lem:inverse-descent}
forces \eqref{marked-local:eq:inverse-T} as the linear central deck action.
The elliptic calculation \eqref{marked-local:eq:actual-elliptic-deck} agrees
with its action on the first complex coordinate. Conversely, the
elliptic calculation alone does not construct the lift to a
four-torus or identify an affine origin there. The distinction
between the linear map and an affine lift will matter below.

We first compute the complete chains for the quotient with deck
matrix $T$. We then use the inverse deck matrix to obtain positive
transport $T$ with the same fibre marking. Both quotients are
explicit geometric spaces; their identification with any prescribed
degeneration requires the corresponding affine descent data.


\section{The affine quotient}

Let $F=(\Z^4\otimes\mathbb R)/\Z^4$, with ordered lattice basis
$(e_1,e_2,e_3,e_4)$. Consider the matrix and translation
\begin{equation}\label{marked-local:eq:native-data}
 T=\begin{pmatrix}
 0&-1&-1&0\\1&-1&0&0\\0&0&1&0\\0&0&0&1
 \end{pmatrix},\qquad a=\tfrac13e_4.
\end{equation}
The affine map $g:F\to F$ is $g(x)=Tx+a$.
Since $T^3=I$, $Ta=a$, and $3a=e_4$ is an integral period,
$g^3=1$. Its first and second powers translate the fourth coordinate
by $1/3$ and $2/3$, so neither has a fixed point. Thus $g$ generates
a free action of the cyclic group of order three.

Write $D=\{s\in\mathbb C:|s|\le1\}$ and
$\zeta=\exp(2\pi i/3)$. The filling considered here is the smooth
quotient
\begin{equation}\label{marked-local:eq:filling}
 N=(F\times D)/\bigl((x,s)\sim(gx,\zeta s)\bigr),
 \qquad q:N\longrightarrow D,\quad q[x,s]=s^3.
\end{equation}
The action on the product is free because its action on $F$ is free.
Local quotient charts are therefore ordinary smooth charts, and
$\partial N=(F\times S^1)/\langle(g,\zeta)\rangle$.
The central fibre has reduced space $G=F/\langle g\rangle$.
In the transverse disk coordinate the map is $s\mapsto s^3$, which
is the meaning of its multiplicity three.

In \eqref{marked-local:eq:filling}, the matrix $T$ is explicitly the action of
the deck generator over $s\mapsto\zeta s$. This convention specifies
a local quotient. It does not identify $T$ with positive boundary
transport. Choosing the translation $a=e_4/3$ fixes its affine action. The freeness follows from this
translation and is used throughout. An order-three translation in a
different invariant direction need not give a free affine action.

Introduce the following integral basis of the first three coordinates:
\begin{equation}\label{marked-local:eq:permutation-basis}
 u_0=e_3,\qquad u_1=-e_1+e_3,\qquad
 u_2=-e_1-e_2+e_3.
\end{equation}
Their column matrix has determinant one. Direct multiplication gives
$Tu_0=u_1$, $Tu_1=u_2$, and $Tu_2=u_0$.
Let $P$ denote this cyclic permutation of the three-torus
$X=(\mathbb R/\Z)^3$, and set $R=P^{-1}$.
If $t\in\mathbb R/\Z$ denotes the $e_4$ coordinate, then
\begin{equation}\label{marked-local:eq:affine-product}
 F=X\times S^1_t,\qquad g(x,t)=(Px,t+1/3).
\end{equation}
The invariant circle $u_0+u_1+u_2$ is
$-2e_1-e_2+3e_3$; thus this basis retains the integral overlattice.

The computations below apply to the actual neighbourhood of a smooth
good reduction with the specified central action. The precise smooth
comparison is as follows.

\begin{proposition}\label{marked-local:prop:family}
Let $\mathcal F\to D$ be a proper smooth family of four-tori,
equipped with an order-three diffeomorphism covering $s\mapsto\zeta s$.
Suppose its action on the central fibre, in the stated marking, is $g$.
Then there is a smooth equivariant trivialization
$\mathcal F\cong F\times D$ whose action is $(g,\zeta)$.
Consequently its quotient neighbourhood is diffeomorphic to $N$.
\end{proposition}
\begin{proof}
Choose a smooth Riemannian metric on the total space and average it
over the finite group. The orthogonal complement to the tangent spaces
of the fibres is a group-invariant horizontal distribution. Lift each
radial path $v\mapsto vs$, $0\le v\le1$, horizontally from the
central fibre. Properness on the closed disk ensures that these lifts
exist for the entire interval. Smooth dependence for the differential
equation gives a smooth trivialization, with inverse obtained by lifting
the reversed radial paths. Rotation sends the radial path to $s$ to
the radial path to $\zeta s$. Invariance of the horizontal distribution
therefore makes the trivialization equivariant. On the central fibre
it is the identity, so the fibre action is exactly $g$.
Taking the free finite quotients gives the claimed diffeomorphism.
\end{proof}

For a family with its torus group structure, there is a useful choice
that also preserves the marked translations. The kernels of the
exponential maps from the real Lie algebras of the fibres form a locally
constant lattice of rank four. Over the disk choose its four marked
periods $v_j(s)$. They form a real basis of each Lie algebra. The map
\[
 (c_1,c_2,c_3,c_4)\longmapsto
       \exp_s\!\left(\sum_jc_jv_j(s)\right)
\]
descends from $\mathbb R^4$ to a smooth group isomorphism
$\mathbb R^4/\mathbb Z^4\to\mathcal F_s$.
The integral matrix of the linear deck action is locally constant,
hence constant on the disk. Translation by $v_4(s)/3$ becomes
translation by $e_4/3$. Thus these period coordinates simultaneously
give \eqref{marked-local:eq:native-data} and its chosen torsion translation.
On the unit boundary circle, the coefficient of $v_4(s)$ in
$(\log s)/(2\pi i)\,v_4(s)$ is real. It therefore gives exactly
the angular translation used below. This fixes the boundary marking
when the smooth group family and its periods are supplied.

This comparison is smooth. It does not assert that a varying complex
torus family is holomorphically a product. The hypotheses state exactly
which part of good reduction is needed for the local homology calculation.

\section{The quotient core and the disk coordinate}

For the homeomorphism $R$ of $X$, its mapping torus is
\[
 G_R=(X\times[0,1])/((x,1)\sim(Rx,0)).
\]
Equivalently, use pairs $(x,\tau)\in X\times\mathbb R$ with
$(x,\tau+1)\sim(Rx,\tau)$. Brackets denote their classes.

\begin{proposition}\label{marked-local:prop:disk-product}
The quotient core $G=F/\langle g\rangle$ is diffeomorphic to $G_R$,
and the map
\begin{equation}\label{marked-local:eq:disk-product}
 \Phi:N\longrightarrow G_R\times D,
 \qquad [x,t,s]\longmapsto\bigl([x,3t],e^{-2\pi it}s\bigr)
\end{equation}
is a diffeomorphism. In these coordinates, writing $z=e^{-2\pi it}s$,
the degeneration map is
\begin{equation}\label{marked-local:eq:degeneration-map}
                   q([x,\tau],z)=e^{2\pi i\tau}z^3.
\end{equation}
\end{proposition}
\begin{proof}
The map $(x,t)\mapsto[x,3t]$ is invariant under $g$, because
\[
 [Px,3t+1]=[RPx,3t]=[x,3t].
\]
It is also invariant under $t\mapsto t+1$, since $R^3=I$.
Every point of $G_R$ has a preimage, and its three preimages in $F$
are precisely one orbit of $g$. The resulting map $G\to G_R$
is a diffeomorphism in the local product charts.

The second coordinate in \eqref{marked-local:eq:disk-product} is unchanged by
$t\mapsto t+1$ and by $(t,s)\mapsto(t+1/3,\zeta s)$.
An inverse is obtained by choosing a representative $[x,\tau]$,
setting $t=\tau/3$, and taking $s=e^{2\pi it}z$.
Changing the representative changes $(x,t,s)$ by the defining finite
action or an integral period. Thus the inverse is well-defined and
smooth. Cubing $s=e^{2\pi it}z$ proves
\eqref{marked-local:eq:degeneration-map}.
\end{proof}

The map $r:N\to G_R$, given by the first coordinate of $\Phi$,
is a deformation retraction after identifying $G_R$ with the zero
section: multiply $z$ by a real number decreasing from one to zero.
The boundary is $G_R\times S^1_z$. Although the disk bundle is smooth
and trivial, the map $q$ has a multiple fibre and a nontrivial boundary
monodromy. Formula~\eqref{marked-local:eq:degeneration-map} retains this distinction.

Orient $F$ by $(u_0,u_1,u_2,t)$, orient the disk by its complex
coordinate, and give $N$ the resulting quotient orientation.
Orient $G_R$ by $(u_0,u_1,u_2,\tau)$. Then $\Phi$ preserves the
orientation of $N$ and $G_R\times D$. The boundary orientation of
$N$ is the product orientation of $G_R$ and the positive circle $S^1_z$.

\section{The boundary marking and its two orientations}

We now identify a smooth fibre and its transported chains inside this
specific filling. Take the boundary value $q=1$. Its fibre is the
image of $s=1$, with marking
\begin{equation}\label{marked-local:eq:fibre-inclusion}
 F\longrightarrow N,\qquad (x,t)\longmapsto[x,t,1].
\end{equation}
After the deformation retraction, this map becomes the covering
\begin{equation}\label{marked-local:eq:cover-map}
             f:F\longrightarrow G_R,
             \qquad f(x,t)=[x,3t].
\end{equation}
Here $f$ denotes this local fibre map, and $q$ remains the map to the disk.

Let $M=P\times I$ act on $F=X\times S^1_t$.
Its mapping torus $B_M$ is formed by
$(x,t,1)\sim(Px,t,0)$, where the last variable is denoted by $v$.
The following boundary map removes the affine translation while keeping
its effect on the filling map:
\begin{equation}\label{marked-local:eq:clockwise-boundary}
 j_-:B_M\longrightarrow\partial N,
 \qquad [x,t,v]\longmapsto
       [x,t-v/3,e^{-2\pi iv/3}].
\end{equation}
Indeed, at $v=1$, applying the quotient generator gives
$[x,t-1/3,\zeta^{-1}]=[Px,t,1]$, as required at the other end.
Moreover $qj_-=e^{-2\pi iv}$. Thus $v$ parametrizes the clockwise
circle in the base disk. The term $-v/3$ is the logarithmic translation
$\sigma(s)=(\log s)/(2\pi i)\,e_4$ along this chosen root of $q$.

In the coordinates of Proposition~\ref{marked-local:prop:disk-product},
\begin{equation}\label{marked-local:eq:boundary-product}
       \Phi j_-([x,t,v])=([x,3t-v],e^{-2\pi it}).
\end{equation}
The map on the two circle coordinates is integral with matrix
$\left(\begin{smallmatrix}3&-1\\-1&0\end{smallmatrix}\right)$,
of determinant $-1$. The mapping-torus end relations give exactly the
corresponding fibre permutations. Solving $t=-\arg(z)/(2\pi)$ and
$v=3t-\tau$ gives its inverse, with changes of lifts absorbed by
these relations. Hence $j_-$ is a diffeomorphism. With the orientation
of $B_M$ given by the fibre followed by increasing $v$, it reverses
the outward boundary orientation of $N$.

For the counterclockwise circle the linear return map is $M^{-1}$,
and the orientation-preserving boundary identification is
\begin{equation}\label{marked-local:eq:counterclockwise-boundary}
 j_+:B_{M^{-1}}\longrightarrow\partial N,
 \quad [x,t,v]\longmapsto[x,t+v/3,e^{2\pi iv/3}],
 \quad \Phi j_+=([x,3t+v],e^{-2\pi it}).
\end{equation}
Its circle-coordinate determinant is $+1$.
Thus the matrix $T$ in \eqref{marked-local:eq:native-data} is the deck action
and the clockwise linear boundary transport in these coordinates.
Inverting the based meridian inverts this transport. Both orientations
will be retained at chain level.

The composite of the clockwise boundary map with the retraction is
\begin{equation}\label{marked-local:eq:actual-prism}
 K_-:B_M\longrightarrow G_R,
 \qquad K_-([x,t,v])=[x,3t-v].
\end{equation}
On the interval before gluing, this is a homotopy from $f$ to $fM$:
$[x,3t-1]=[Px,3t]$. This explicit homotopy will determine the
prism map on chains.

\section{Finite integral chains}

A chain complex consists of abelian groups with homomorphisms
$d_n:C_n\to C_{n-1}$ satisfying $d_{n-1}d_n=0$. A chain map
commutes with these boundaries. A chain homotopy from $A$ to $B$
is a degree-one map $H$ with $dH+Hd=B-A$.
The homology in degree $n$ is $\ker d_n/\im d_{n+1}$.
We use free groups generated by oriented cells, with the boundary
given by their attaching maps. Interval subdivisions and their prisms
compare these chains with singular chains, the finite integer sums
of continuous oriented simplices with alternating face boundary.

Give each circle in $X$ one vertex and one oriented edge. The products
give one cell for every subset of the three coordinates. Opposite
faces cancel, so the cellular chain groups have zero differential.
Write $D_k$ for the group in degree $k$ and $D_k=0$ outside $0\le k\le3$.
The bases are
\begin{equation}\label{marked-local:eq:cell-bases}
 \begin{array}{c|c}
 k&\text{basis of }D_k\\ \hline
 0&1\\
 1&u_0,u_1,u_2\\
 2&w_0=u_1\wedge u_2,
     \quad w_1=u_2\wedge u_0,
     \quad w_2=u_0\wedge u_1\\
 3&\omega=u_0\wedge u_1\wedge u_2.
 \end{array}
\end{equation}
The coordinate permutation is an actual cellular homeomorphism. It acts
by the identity on $D_0,D_3$ and cyclically on both displayed triples.
Denote these chain actions by $P_k$, put $R_k=P_k^{-1}$, and set
\begin{equation}\label{marked-local:eq:norm}
                 S_k=I+R_k+R_k^2.
\end{equation}
Thus $S_0=S_3=3$, whereas $S_1$ and $S_2$ are the three-by-three
matrix all of whose entries are one.

For products with an interval, orient its chain before the fibre chain.
If $z$ is a fibre chain, write $\sigma z=[0,1]\times z$.
Its boundary is the terminal chain minus the initial chain minus
$\sigma\partial z$. In particular, the orientation of $\sigma z$
is $(-1)^k$ times the fibre-first orientation when $z$ has degree $k$.

The fibre $F=X\times S^1_t$ has the finite chain complex
\begin{equation}\label{marked-local:eq:fibre-complex}
 C_n=D_n\oplus D_{n-1},\qquad d_C=0,
 \qquad M_n=P_n\oplus P_{n-1}.
\end{equation}
The second component is the oriented $t$-circle prism.
For the core $G_R$, hence also the deformation retract of $N$, take
\begin{equation}\label{marked-local:eq:filling-complex}
 Q_n=D_n\oplus D_{n-1},\qquad
 d_Q(x,y)=((R_{n-1}-I)y,0).
\end{equation}
This follows directly by attaching one $\tau$-prism to each cell of $X$.
The ranks are $1,4,6,4,1$. The only nonzero differentials are
$d_{Q,2}$ and $d_{Q,3}$, acting by $R-I$ on the suspended triples.
Thus \eqref{marked-local:eq:filling-complex} specifies every boundary through degree
four, with its bases and integer entries.

The clockwise boundary mapping torus has complex
\begin{equation}\label{marked-local:eq:boundary-complex}
 B_n=C_n\oplus C_{n-1},\qquad
 d_B(a,b)=((M_{n-1}-I)b,0).
\end{equation}
This is its actual cellular complex because $M$ permutes the product
cells. No simultaneous replacement of singular chains by exterior
homology is needed: these specific cellular actions exist on the stated
spaces and give their interval attachments directly.

\begin{theorem}\label{marked-local:thm:maps}
For the fibre inclusion and the clockwise boundary map above, the finite
chain maps and prism homotopy are
\begin{equation}\label{marked-local:eq:F-H}
 F_n(x,y)=(x,S_{n-1}y),\qquad
 H_n(x,y)=(0,-P_nx)\in Q_{n+1}.
\end{equation}
They satisfy
\begin{equation}\label{marked-local:eq:homotopy-identity}
 d_QF=0=Fd_C,\qquad d_QH+Hd_C=F(M-I).
\end{equation}
In degree $n$, write a boundary chain as
$(x,y,z,w)$ in
$D_n\oplus D_{n-1}\oplus D_{n-1}\oplus D_{n-2}$.
The actual boundary inclusion followed by retraction is represented by
\begin{equation}\label{marked-local:eq:K}
 K_n(x,y,z,w)=(x,S_{n-1}y-P_{n-1}z).
\end{equation}
These formulas hold in every degree, including the entire range needed
for gluing boundaries in degrees two, three and four.
\end{theorem}
\begin{proof}
The map $f(x,t)=[x,3t]$ sends a cell at $t=0$ identically to the
corresponding core fibre cell. A $t$-prism traverses three successive
$\tau$-intervals. On these intervals its fibre chains are transported
by $I,R,R^2$, respectively. Subdivide the source interval into three
equal parts and add their oriented images. This gives $S_{n-1}$,
and proves the formula for $F$.

The homotopy \eqref{marked-local:eq:actual-prism} moves a fibre cell at $t=0$
backwards through one core interval. A backward interval has chain
$-\sigma(Px)$, because its terminal chain is $Px$ and its initial
chain is $x$. Indeed,
$(R-I)(-P)x=(P-I)x$. On a $t$-prism over a $k$-cell, the
homotopy factors through that cell and the one-dimensional core circle:
its circle coordinate is $3t-v$. Subdivide the square in $(t,v)$
along the finitely many integer levels of $3t-v$. In every piece its
image lies in the product of the same $k$-cell, possibly permuted,
and an interval. It lies in the $(k+1)$-skeleton, so its cellular
component in degree $k+2$ is zero. This proves the formula for $H$
on the second summand as well.

All subdivisions are compatible with coordinate faces. Their interval
and square prisms give chain comparisons with singular chains and
with the geometric maps. In particular, the displayed formulas come
from the stated covering and boundary homotopy, with their common
cell markings.

For an algebraic verification, $(R-I)S=R^3-I=0$ and
$S(P-I)=0$, while $(R-I)(-P)=P-I$. Consequently
\[
 d_QH(x,y)=((P_n-I)x,0)=F(M-I)(x,y),
\]
which proves \eqref{marked-local:eq:homotopy-identity}. Combining the fibre part
$F$ and the interval part $H$ gives \eqref{marked-local:eq:K}.
\end{proof}

For the counterclockwise boundary the corresponding homotopy is
\begin{equation}\label{marked-local:eq:positive-H}
 H_n^+(x,y)=(0,x),\qquad d_QH^+=F(M^{-1}-I).
\end{equation}
Its geometric representative is $[x,3t+v]$ from
\eqref{marked-local:eq:counterclockwise-boundary}.
The reversal map from the clockwise boundary complex to the
counterclockwise one is
\[
 J(a,b)=(a,-Mb).
\]
Indeed, $(M^{-1}-I)(-M)=M-I$, so it commutes with boundaries.
Moreover $F+H^+(-M)=F+H$, hence $K_+J=K_-$.
These equations give the explicit comparison when a gluing uses the
opposite meridian convention.

\section{Homology and specialization}

Since the boundary in
\eqref{marked-local:eq:filling-complex} acts only on the suspended summand,
\begin{equation}\label{marked-local:eq:Q-homology}
 H_n(Q)=\coker(R_n-I)\oplus\ker(R_{n-1}-I).
\end{equation}
For a cyclic permutation of three generators, the image of $R-I$
is exactly the lattice of vectors whose coordinates sum to zero.
The differences of two generators span that lattice: if
$a_0+a_1+a_2=0$, then
$a_0u_0+a_1u_1+a_2u_2=a_0(u_0-u_2)+a_1(u_1-u_2)$.
The kernel is generated by their sum. Thus the cokernel and kernel
are both infinite cyclic and have the following integral generators:
\[
 \begin{array}{c|cc}
 k&\coker(R_k-I)&\ker(R_k-I)\\ \hline
 0&1&1\\
 1&[u_0]&U=u_0+u_1+u_2\\
 2&[w_0]&W=w_0+w_1+w_2\\
 3&\omega&\omega.
 \end{array}
\]
Each displayed vector generates its full integral kernel or cokernel.
No quotient here has torsion.

\begin{theorem}\label{marked-local:thm:homology}
The integral homology of $N$ is
\[
 H_0(N)=\Z,\quad H_1(N)=\Z^2,\quad H_2(N)=\Z^2,
 \quad H_3(N)=\Z^2,\quad H_4(N)=\Z,
\]
and zero in higher degrees. In the ordered bases of
\eqref{marked-local:eq:cell-bases} and \eqref{marked-local:eq:Q-homology}, the fibre map on
homology is
\begin{equation}\label{marked-local:eq:F-homology}
\begin{array}{c|c|c|c}
 n&F_*&\ker F_*&\coker F_*\\ \hline
 0&1&0&0\\
 1&(a_0,a_1,a_2;b)\mapsto(a_0+a_1+a_2,3b)&\Z^2&\Z/3\Z\\
 2&(a_0,a_1,a_2;b_0,b_1,b_2)\mapsto
       (a_0+a_1+a_2,b_0+b_1+b_2)&\Z^4&0\\
 3&(a;b_0,b_1,b_2)\mapsto(a,b_0+b_1+b_2)&\Z^2&0\\
 4&a\mapsto3a&0&\Z/3\Z.
\end{array}
\end{equation}
Each kernel in degrees one, two and three is the sum-zero lattice
in the indicated triple or pair of triples.
\end{theorem}
\begin{proof}
The preceding image and kernel computation applied to
\eqref{marked-local:eq:Q-homology} gives all the homology groups.
In degrees one and two, a vector in a triple maps to its coordinate
sum in the cokernel, while $S$ sends it to that sum times $U$ or $W$.
On $D_0,D_3$, $S=3$. Applying these facts to the two summands
of $F_n$ in \eqref{marked-local:eq:F-H} gives every row of
\eqref{marked-local:eq:F-homology}, together with the displayed kernels and
cokernels. The top-dimensional map has degree three, as does the
covering $f:F\to G_R$.
\end{proof}

For cohomology use the dual cellular complex
$\operatorname{Hom}(Q_*,\Z)$, whose coboundary is the transpose of
the chain boundary. The kernel and image splittings just given show
that evaluation identifies its cohomology with the integral dual of
\eqref{marked-local:eq:Q-homology}. Let $\eta\in H^1(F;\Z)$ evaluate to one
on the oriented $t$ circle and to zero on all $u_i$.
Let $u_i^*,w_i^*,\omega^*$ denote
the dual fibre generators. Products with $\eta$ in the following
description use the order of the suspended summand of $C$.

The invariant fibre cohomology has generators
\[
\begin{array}{c|c}
 1&u_0^*+u_1^*+u_2^*,\quad\eta\\
 2&w_0^*+w_1^*+w_2^*,\quad
       \eta\wedge(u_0^*+u_1^*+u_2^*)\\
 3&\omega^*,\quad\eta\wedge(w_0^*+w_1^*+w_2^*)\\
 4&\eta\wedge\omega^*.
\end{array}
\]
Here the exterior forms are identified with the dual product cells.
In degree two, for example, the signed cyclic basis for the $w_i$
is the same basis as in \eqref{marked-local:eq:cell-bases}; no change to an
unsigned lexicographic basis is implicit.

The pullback $F^*:H^n(N;\Z)\to H^n(F;\Z)^M$ is the
cohomological specialization map for this neighbourhood.
Taking the transpose of \eqref{marked-local:eq:F-homology} shows that it is
injective. Its image is the
full invariant lattice in degrees two and three. In degree one the
second generator is multiplied by three, and in degree four its sole
generator is multiplied by three. Consequently
\begin{equation}\label{marked-local:eq:specialization-indices}
 [H^n(F;\Z)^M:\im F^*]=3,1,1,3
                         \quad\text{for }n=1,2,3,4.
\end{equation}
In the original marking,
$u_0^*+u_1^*+u_2^*=e_3^*$ and $\eta=e_4^*$.
Thus degree-one specialization is precisely
$\Z e_3^*\oplus3\Z e_4^*$ inside
$\Z e_3^*\oplus\Z e_4^*$.

\section{The full boundary inclusion}

The disk trivialization makes the boundary inclusion followed by
retraction the projection $G_R\times S^1\to G_R$.
Its complete integral content can also be read directly from
\eqref{marked-local:eq:K}.

\begin{proposition}\label{marked-local:prop:boundary-splitting}
The chain map $K:B\to Q$ is surjective and has a chain section
\begin{equation}\label{marked-local:eq:section}
 s_n(x,b)=(x,0,-R_{n-1}b,0).
\end{equation}
Its kernel is isomorphic to the shifted chain complex $Q[1]$, where
$Q[1]_n=Q_{n-1}$ and its differential is $-d_Q$.
Consequently
\[
 H_n(\partial N;\Z)=H_n(N;\Z)\oplus H_{n-1}(N;\Z),
 \qquad K_*\text{ is projection onto the first summand}.
\]
In particular the boundary homology is torsion-free, with ranks
$1,3,4,4,3,1$ in degrees zero through five.
\end{proposition}
\begin{proof}
Equation~\eqref{marked-local:eq:K} gives $Ks=I$. The boundary of the suspended
part in \eqref{marked-local:eq:section} is
$(P-I)(-R)b=(R-I)b$, proving $d_Bs=sd_Q$.
A vector in $\ker K_n$ has $x=0$ and
$z=P_{n-1}^{-1}S_{n-1}y=S_{n-1}y$. Thus it is uniquely of the
form $(0,y,Sy,w)$, with differential
\[
                    (y,w)\longmapsto((P-I)w,0).
\]
The map from $Q[1]_n=D_{n-1}\oplus D_{n-2}$ sending
$(a,b)$ to $(P_{n-1}a,b)$ in these kernel coordinates is invertible.
It commutes with differentials because $-P(R-I)=P-I$.
Therefore $B=s(Q)\oplus\ker K$ is the asserted direct sum of
chain complexes. Taking homology gives the formula and all six ranks.
The top group maps to zero, since $H_5(N)=0$.
\end{proof}

The section has a geometric representative: in
$\partial N=G_R\times S^1_z$, fix $z=1$. Its core interval
traverses the clockwise boundary interval backwards, yielding
$-R$ in \eqref{marked-local:eq:section}. The kernel describes the extra normal
circle. The chain identification with $Q[1]$ fixes its signs by the
displayed formulas, rather than by an unnamed choice of splitting.

This determines the specified local filling, its entire finite chain
complex, and its actual fibre and boundary maps. It uses the particular affine
translation $e_4/3$, the free action in \eqref{marked-local:eq:affine-product},
and the smooth good-reduction comparison of
Proposition~\ref{marked-local:prop:family}. Changing the translation requires a
new quotient and comparison. The other finite-order neighbourhood and
the cusp degeneration require their own chain attachments. The local
cell models must also be compared compatibly with a common complex
carrying all the global transport maps. The local
calculation alone gives no homology calculation for their global gluing
and no assertion about sphere topology or a complex structure on a sphere.
\section{The quotient with the prescribed positive transport}

Keep the lattice, the oriented basis $(u_0,u_1,u_2,e_4)$, and the
cell groups $D_k$ already defined. Now take $R=P^{-1}$ as the deck
action on $X$, with the same translation $a=e_4/3$. Define
\begin{equation}\label{marked-local:eq:positive-quotient}
 N^+=(X\times S^1_t\times D)/
 \bigl((x,t,s)\sim(Rx,t+1/3,\zeta s)\bigr),
 \qquad q[x,t,s]=s^3.
\end{equation}
This action is free by its fourth coordinate, just as for $N$.
Here $t$ is again the fibre circle coordinate; $q$ denotes the
coordinate on the quotient disk. Let $G_P$ be the mapping torus
with relation $[x,\tau+1]=[Px,\tau]$.

\begin{theorem}\label{marked-local:thm:positive-geometry}
There is an orientation-preserving diffeomorphism
\begin{equation}\label{marked-local:eq:positive-product}
 \Phi^+:N^+\longrightarrow G_P\times D,
 \qquad [x,t,s]\longmapsto([x,3t],e^{-2\pi it}s).
\end{equation}
In these coordinates $q=e^{2\pi i\tau}z^3$.
The positively oriented boundary, with fibre marking at $q=1$,
is identified with $B_M$, where $M=P\times I$, by
\begin{equation}\label{marked-local:eq:positive-native-boundary}
 j:B_M\longrightarrow\partial N^+,
 \qquad [x,t,v]\longmapsto[x,t+v/3,e^{2\pi iv/3}].
\end{equation}
The fibre map and the boundary map followed by retraction are
\begin{equation}\label{marked-local:eq:positive-geometric-maps}
 f^+(x,t)=[x,3t],\qquad
 k^+([x,t,v])=[x,3t+v].
\end{equation}
Thus the positive transport in this marking is exactly $T$.
\end{theorem}
\begin{proof}
Under the quotient generator,
$[Rx,3t+1]=[PRx,3t]=[x,3t]$, and the factors
$e^{-2\pi i/3}$ and $\zeta$ cancel in the disk coordinate.
The formula is therefore well defined. A lift of $\tau$ gives
$t=\tau/3$ and $s=e^{2\pi i\tau/3}z$; changing the lift is exactly
the quotient relation. This proves that the formula and its inverse
are smooth, including at $z=0$. Its real Jacobian has positive
factor $3$ in the $t$ direction and a rotation in the disk.
The formula for $q$ follows by cubing $s$.

At $v=1$ in \eqref{marked-local:eq:positive-native-boundary}, applying the inverse
deck generator gives $[Px,t,1]$. This is the endpoint identification
of $B_M$, and $qj=e^{2\pi iv}$. Under $\Phi^+$ the map is
$([x,3t+v],e^{-2\pi it})$. The two circle coordinates have matrix
$\left(\begin{smallmatrix}3&1\\-1&0\end{smallmatrix}\right)$
of determinant one. Solving these coordinates as before gives a
smooth inverse compatible with the mapping-torus relations.
The boundary orientation is therefore preserved. Setting the disk
coordinate equal to zero gives both maps in
\eqref{marked-local:eq:positive-geometric-maps}.
\end{proof}

No reversal of the positive base circle was used. For comparison,
the map $\mathcal J:B_M\to B_{M^{-1}}$ given by
$\mathcal J([x,t,v])=[Px,t,1-v]$ satisfies
$j_+\mathcal J=j_-$ for the earlier quotient $N$. It reverses the
base circle. There cannot be a comparison between these two earlier
markings covering the identity on the oriented base and inducing
the identity on fibre homology: it would intertwine $P$ and $R$
by the identity. But $Pu_0=u_1$ and $Ru_0=u_2$.

\section{Integral maps for positive transport}

The cell model for the core of $N^+$ is
\begin{equation}\label{marked-local:eq:Q-positive}
 Q^+_n=D_n\oplus D_{n-1},\qquad
 d_+(x,y)=((P_{n-1}-I)y,0).
\end{equation}
The fibre complex remains $C$ and its boundary mapping torus remains
$B$, with differential $(a,b)\mapsto((M-I)b,0)$.
Set $S=I+P+P^2=I+R+R^2$ in each exterior degree.

\begin{theorem}\label{marked-local:thm:positive-chains}
With the interval-first cell orientation, the geometric maps
\eqref{marked-local:eq:positive-geometric-maps} induce
\begin{equation}\label{marked-local:eq:positive-FHK}
 \begin{split}
 F^+_n(x,y)&=(x,S_{n-1}y),\\
 H^+_n(x,y)&=(0,x),\\
 K^+_n(x,y,z,w)&=(x,S_{n-1}y+z).
 \end{split}
\end{equation}
They satisfy
\begin{equation}\label{marked-local:eq:positive-identities}
 d_+F^+=0,\qquad d_+H^++H^+d_C=F^+(M-I),
 \qquad d_+K^+=K^+d_B.
\end{equation}
These are maps on the actual fibre, core and boundary cell complexes,
not merely maps on their homology.
\end{theorem}
\begin{proof}
The cover $f^+$ traverses three core intervals in the positive
direction. The successive transports are $I,P,P^2$, which give
$S$ on the suspended summand and the identity on fibre cells.
The homotopy $[x,3t+v]$ traverses one positive core interval on
a cell at $t=0$. Its chain is $\sigma x$, giving $H^+(x,0)=(0,x)$.
On the $t$-prism of a $k$-cell, subdivide the $(t,v)$ square along
integer values of $3t+v$. Each image lies in the product of a
permuted $k$-cell and a core interval, hence in the $(k+1)$-skeleton.
Its degree-$(k+2)$ cellular component is zero. This proves the second
formula on every summand. The subdivisions agree on faces, so these
are cellular representatives of the covering and its stated homotopy.
Combining their fibre and boundary-interval components gives $K^+$.

For the chain identities, $(P-I)S=S(P-I)=0$ and
$d_+H^+(x,y)=((P-I)x,0)=F^+(M-I)(x,y)$.
The equality for $K^+$ follows by applying these two identities to
the fibre and suspended boundary summands separately.
\end{proof}

All entries through degree four can be read without an implicit
exterior operation. The groups $D_0,D_3$ have rank one and
$P_0=P_3=1$; $D_1,D_2$ have rank three and
\[
 P_1=P_2=\begin{pmatrix}0&0&1\\1&0&0\\0&1&0\end{pmatrix},
 \qquad S_0=S_3=3,\qquad
 S_1=S_2=\begin{pmatrix}1&1&1\\1&1&1\\1&1&1\end{pmatrix}.
\]
Thus $Q^+$ has ranks $1,4,6,4,1$, with only $d_{+,2}$ and
$d_{+,3}$ nonzero, each using the displayed $P-I$ on its suspended
triple. In particular $F^+_4$ is multiplication by three, $H^+_3$
sends the top fibre cell $\omega$ to its core prism, and $H^+_4=0$.
These formulas specify the maps in every degree.

\begin{proposition}\label{marked-local:prop:positive-splitting}
The positive boundary map has chain section
$s^+_n(x,b)=(x,0,b,0)$. Its kernel consists exactly of
$(0,y,-Sy,w)$, with differential $(y,w)\mapsto((P-I)w,0)$.
The map $(a,b)\mapsto(-a,b)$ identifies $Q^+[1]$ with this kernel.
Consequently $B\cong Q^+\oplus Q^+[1]$ as integral chain complexes.
\end{proposition}
\begin{proof}
The formula for $K^+$ gives $K^+s^+=I$. Since the boundary of
its third component is $(P-I)b$, $s^+$ commutes with boundaries.
The equation $K^+=0$ forces $x=0,z=-Sy$. Its boundary has
second component $(P-I)w$ and third component zero, because
$(P-I)S=0$. This proves the kernel formula. On $Q^+[1]$ the
differential is $-d_+$; multiplication by $-1$ on its first
summand changes this to the displayed kernel differential.
All changes of basis are integral and invertible.
\end{proof}

The section is the actual normal-circle section $z=1$ in
$G_P\times S^1_z$: a positive core interval traverses the positive
boundary interval, with no inverse transport. The kernel carries
the extra normal circle. The image of $P-I$, like that of $R-I$,
is the full sum-zero lattice in each cyclic triple, and its kernel
is generated by the sum. Hence the proof of
\eqref{marked-local:eq:Q-homology} applies with $P$ in place of $R$: explicitly
$H_n(Q^+)=\coker(P_n-I)\oplus\ker(P_{n-1}-I)$, with the same
integral generators $1,[u_0],[w_0],\omega$ and $1,U,W,\omega$.
The covering uses exactly the same $S$. Every homology matrix in
\eqref{marked-local:eq:F-homology} therefore gives $F^+_*$ in these stated bases.
Its kernels are the indicated sum-zero lattices, and its only
nonzero cokernels are $\mathbb Z/3\mathbb Z$ in degrees one and four.
The groups of $N^+$ are torsion-free, of ranks $1,2,2,2,1$.
The boundary groups are torsion-free, of ranks $1,3,4,4,3,1$,
and $K^+_*$ is the split projection. Dualizing these free groups
gives specialization indices $3,1,1,3$ and degree-one image
$\mathbb Ze_3^*\oplus3\mathbb Ze_4^*$ in the unchanged fibre marking.

\section{The affine comparison required by a prescribed filling}

A marking of homology does not specify an origin of a torus.
In period coordinates an affine transformation has the form
$x\mapsto Ax+b$, where $b$ is a point of the torus. Changing the
origin by $c$ changes its translation to $b+(A-I)c$.
It does not change its homology matrix. A comparison of fillings
with a chosen boundary section must retain that section as well
as the integral periods.

\begin{proposition}\label{marked-local:prop:marked-comparison}
Let a proper smooth torus group family over the $s$ disk have
the marked periods \eqref{marked-local:eq:period-lattice}. Suppose its original
deck transformation over $s\mapsto\zeta s$ preserves its zero
section, has matrix $A=T^{-1}$, and the chosen origin on the
punctured quotient is the descent of that zero section. Twist the
deck transformation by the section $e_4(s)/3$. Then its quotient
is smoothly identified with $N^+$, with the same positive fibre
marking and boundary map \eqref{marked-local:eq:positive-native-boundary}.
In particular \eqref{marked-local:eq:Q-positive} and
\eqref{marked-local:eq:positive-FHK} give its integral filling data.
\end{proposition}
\begin{proof}
The periods trivialize the real Lie algebras and exponentiate to
smooth group identifications with $\mathbb R^4/\mathbb Z^4$.
The deck matrix is locally constant, hence $A$ throughout the disk.
Preservation of zero makes its affine translation zero. The twisted
action is therefore $(x,s)\mapsto(Ax+e_4/3,\zeta s)$.
In the oriented $u$ basis this is \eqref{marked-local:eq:positive-quotient}.
On the lifted positive boundary path $s=e^{2\pi iv/3}$,
the logarithmic translation
$(\log s)/(2\pi i)\,e_4$ is $ve_4/3$. It vanishes at $v=0$
and gives exactly \eqref{marked-local:eq:positive-native-boundary}.
All subsequent geometric and chain comparisons have already been
constructed for this quotient. No reversal or change of periods is
involved.
\end{proof}

The zero-section hypothesis is substantive. To see its effect while
retaining the complex torus \eqref{marked-local:eq:period-lattice}, let
$b\in\{0,1,2\}$ and define three holomorphic deck actions
\begin{equation}\label{marked-local:eq:affine-ambiguity}
 \rho_b(x,s)=(Ax+be_4/3,\zeta s).
\end{equation}
All have order three, all have the same linear action on the full
integral lattice, and all induce the same action on the elliptic
coordinate. The descended punctured families are isomorphic as
marked torus families. In fact, on the universal logarithm cover put
\[
 \sigma_b(s)=\frac{b\log s}{2\pi i}\,e_4\in\mathbb C^2/\Lambda.
\]
A change of logarithm adds the integral period $be_4$, so this is a
single-valued holomorphic torus section on the punctured disk.
Since $Ae_4=e_4$, it satisfies
$\sigma_b(\zeta s)=A\sigma_b(s)+be_4/3$.
Translation by $\sigma_b$ consequently intertwines $\rho_0$ and
$\rho_b$, induces the identity on homology, and at $s=1$ fixes the
marked fibre point. It carries the descended zero section to a
section of the other punctured quotient.

Nevertheless twisting each action by the same specified $a=e_4/3$
produces translation $(b+1)e_4/3$. The cases $b=0,1$ are free,
whereas $b=2$ has fixed points on the central fibre: the translation
is then an integral period and the origin is fixed by $A$.
Thus positive integral monodromy, the elliptic deck calculation,
the increment $a$, and even an identified punctured family do not
imply freeness or select an affine filling. The examples establish
this precise failure of inference; they do not identify any one
of these alternative descents with a separately prescribed geometric
construction.

There is a further ambiguity even when the affine central action is
fixed. For an integer $r$, the formula
\[
 L_r([x,t,v])=[x,t+rv,v]
\]
defines a diffeomorphism of $B_M$ covering the identity on the
positive base and restricting to the identity on its marked fibre
at $v=0$. The endpoint condition holds because $r$ is integral
and $M$ fixes the $t$ circle. However the two boundary maps $j$
and $jL_r$ have distinct filling effects. Their retractions are
$[x,3t+v]$ and $[x,3t+(3r+1)v]$.
On the loop with $x=t=0$, they give respectively $[\tau]$ and
$(3r+1)[\tau]$ in $H_1(G_P;\mathbb Z)$, where $[\tau]$ has
infinite order. Thus an unspecified change of section along the
meridian changes a genuine attachment map while preserving the
linear monodromy and the marked starting fibre.


\section{The origin section in the line-bundle reduction}

We now specify the local geometric hypotheses under which the
four-torus filling is obtained from \eqref{marked-local:eq:weierstrass}.
Let $O$ be its origin section, let $P_0$ be a section disjoint
from $O$ near $t=0$, and put $\mathcal M=\mathcal O(P_0-O)$.
Write $\mathcal M^\times$ for the complement of its zero section.
Suppose a specified lifted translation acts on $\mathcal M^\times$
and its quotient over the punctured disk is a family of complex
two-tori. Choose the origin of these tori using a nonzero section
$e$ of $\mathcal M|_O$ over the disk.

After $t=s^3$, let $P'_0,O'$ be the transformed sections on
\eqref{marked-local:eq:good-equation}, and put
$\mathcal M'=\mathcal O(P'_0-O')$. We assume the following smooth
reduction data: $P'_0$ is disjoint from $O'$ near zero; the specified
lifted translation extends to $(\mathcal M')^\times$ and commutes
with the natural deck action; its quotient is a proper smooth
family $\mathcal F$ of complex two-tori over the $s$ disk. The
period marking is \eqref{marked-local:eq:period-lattice}, with $\delta$ the
positive circle of the multiplicative fibre and $e_4$ the stated
second vertical period. The positive monodromy is
\eqref{marked-local:eq:positive-T}.

These are hypotheses on the local smooth reduction and its lifted
translation. The calculation below constructs their marked filling
and attachment maps; it does not establish the existence of that
reduction or of the lifted translation from an arbitrary elliptic
family. In particular it does not replace the needed commutation
or properness by a lattice calculation.

\begin{lemma}\label{marked-local:lem:conormal-comparison}
There is a nonzero section $e'$ of $\mathcal M'|_{O'}$ such that
the canonical divisor identification over the punctured disk sends
the pulled-back section $e$ to $s^{-2}e'$. The natural deck action
sends $e'(s)$ to $\zeta e'(\zeta s)$.
\end{lemma}
\begin{proof}
Since $P_0$ does not meet $O$ in this neighbourhood,
$\mathcal O(P_0-O)$ is canonically the ideal line bundle of $O$
there. Restricting this ideal to $O$ gives its conormal line:
the ideal modulo its square. Similarly $\mathcal M'|_{O'}$ is
the conormal line of $O'$, because $P'_0$ is disjoint from it.

The local parameters transverse to the origin sections are
$w=-x/y$ and $W=-X/Y$. To check this at infinity, use the
projective chart with $Y\ne0$: the cubic equation solves $Z/Y$
as a holomorphic function of $X/Y$ near the origin, so $-X/Y$
is a local parameter. The Weierstrass substitution gives the
exact equality
\begin{equation}\label{marked-local:eq:conormal-coordinate}
                         w=s^{-2}W.
\end{equation}
The nonzero section $e$ can therefore be written as $u(t)w$ in
the conormal line, for a holomorphic unit $u$ on a sufficiently
small disk. Define $e'(s)=u(s^3)W$. This is holomorphic and
nonzero, and the canonical divisor identification on $s\ne0$
gives $e=s^{-2}e'$. No choice of a comparison multiplied by an
extra power of $s$ is made here: the parameters are identified
by the actual birational coordinate substitution.

The original pulled-back section is equivariant for the deck
action. If $\rho e'(s)=\lambda(s)e'(\zeta s)$, applying $\rho$
to $e=s^{-2}e'$ consequently yields
$\lambda(s)=(\zeta s)^{-2}/s^{-2}=\zeta$.
The equality extends to zero. Equivalently, the forward action on
the conormal line is the inverse of the normal-coordinate factor
$\zeta^2$ in \eqref{marked-local:eq:actual-elliptic-deck}.
\end{proof}

Use the image of $e'$ as the origin of the smooth torus family
$\mathcal F$. An automorphism between compact complex tori is
affine in these origins: a lift to their complex vector spaces
has lattice-periodic derivative, whose entries are holomorphic
functions on a compact torus and hence constant. Subtracting its
linear part leaves a constant translation. Its linear part is
determined by its map on periods. Lemma~\ref{marked-local:lem:inverse-descent}
therefore gives linear part $A=T^{-1}$, while
Lemma~\ref{marked-local:lem:conormal-comparison} gives translation $\delta/3$.
Indeed multiplication by $e^{2\pi i/3}$ in the multiplicative
fibre is precisely one third of its positive period $\delta$.

Period coordinates make these maps constant over the smooth disk:
their integral linear matrices are locally constant, and the
translation is the specified torsion section $\delta(s)/3$.
Consequently the original deck transformation and its twist by
the unchanged $a_0=e_4/3$ are, respectively,
\begin{equation}\label{marked-local:eq:actual-affine-actions}
 (z,s)\longmapsto(Az+\delta/3,\zeta s),\qquad
 (z,s)\longmapsto(Az+(\delta+e_4)/3,\zeta s).
\end{equation}
Thus the central affine term does not vanish in the origin $e'$.
The second action is free: in the original $e_4$ coordinate its
first two powers translate by $1/3$ and $2/3$.

Along a positive root $s=e^{2\pi iv/3}$ the original origin $e$
differs from $e'$ by $-2v\delta/3$. The logarithmic twist adds
$ve_4/3$. A small nonunit boundary radius can be absorbed by a
constant rescaling of the chosen frame; the angular formula is
unchanged. At the initial point $s=1$ the two origins coincide.
This determines the marked boundary path, including its original
section, and not only the homology monodromy.

\section{The marked local filling and its integral chains}

Write $\delta=U=u_0+u_1+u_2$ and retain the original coordinates
$(x,t)$ in the basis $(u_0,u_1,u_2,e_4)$. Under the preceding
smooth reduction hypotheses, the twisted local filling is
\begin{equation}\label{marked-local:eq:marked-native-quotient}
 N_0=(X\times S^1_t\times D)/
 \bigl((x,t,s)\sim(Rx+U/3,t+1/3,\zeta s)\bigr).
\end{equation}
Its boundary is marked by the original section $e$ and the
counterclockwise meridian. The formula just derived is
\begin{equation}\label{marked-local:eq:marked-native-j}
 j_0([x,t,v])=[x-2vU/3,t+v/3,e^{2\pi iv/3}].
\end{equation}
At $v=1$, the inverse of the generator in
\eqref{marked-local:eq:marked-native-quotient} sends the endpoint to
$[Px-U,t,1]=[Px,t,1]$. Thus its return is $P$, as required in
the original period marking; $qj_0=e^{2\pi iv}$.

\begin{theorem}\label{marked-local:thm:marked-native-filling}
The filling \eqref{marked-local:eq:marked-native-quotient} has the following
explicit orientation-preserving disk trivialization:
\begin{equation}\label{marked-local:eq:marked-native-Phi}
 \Phi_0[x,t,s]=([x+2tU,3t],e^{-2\pi it}s)
                 \in G_P\times D.
\end{equation}
Here $q=e^{2\pi i\tau}z^3$. The marked fibre and positive
boundary maps followed by retraction are
\begin{equation}\label{marked-local:eq:marked-native-FK}
 f_0(x,t)=[x+2tU,3t],\qquad
 k_0([x,t,v])=[x+2tU,3t+v].
\end{equation}
The integral fibre basis
\begin{equation}\label{marked-local:eq:marked-cell-basis}
             u_0,u_1,u_2,v_4,\qquad v_4=e_4-2\delta
\end{equation}
has determinant one relative to $(e_1,e_2,e_3,e_4)$ and is
preserved by $T$, which acts as $P\times I$. In its actual
product cells the filling complex and maps are
\begin{equation}\label{marked-local:eq:marked-final-chains}
 \begin{aligned}
 Q_{0,n}&=D_n\oplus D_{n-1},&d_0(x,y)&=((P-I)y,0),\\
 F_{0,n}(x,y)&=(x,Sy),&H_{0,n}(x,y)&=(0,x),\\
 K_{0,n}(x,y,z,w)&=(x,Sy+z).&&
 \end{aligned}
\end{equation}
In particular $d_0H_0+H_0d_C=F_0(M-I)$ with the positive
native monodromy and the actual marked boundary section.
\end{theorem}
\begin{proof}
Set $y=x+2tU$. The change $(x,t)\mapsto(y,t)$ is an integral
torus diffeomorphism of determinant one, with inverse
$x=y-2tU$. Under the generator of \eqref{marked-local:eq:marked-native-quotient},
\[
 y\longmapsto Rx+U/3+2(t+1/3)U=Ry+U=Ry
                         \quad\text{in }X.
\]
It therefore identifies this quotient with
\eqref{marked-local:eq:positive-quotient} in the $y$ coordinates, without
changing $s$ or the orientation of its disk. It also sends the
marked boundary path to
$[y,t+v/3,e^{2\pi iv/3}]$, since the terms $-2vU/3$ and
$2(v/3)U$ cancel. Applying Theorem~\ref{marked-local:thm:positive-geometry}
gives \eqref{marked-local:eq:marked-native-Phi} and
\eqref{marked-local:eq:marked-native-FK}, with their orientations and inverses.

In the original coordinates, the $t$ loop with $y=0$ has vector
$e_4-2U$. Hence its product-cell basis is exactly
\eqref{marked-local:eq:marked-cell-basis}. Its matrix of columns is
\begin{equation}\label{marked-local:eq:marked-basis-matrix}
 V=\begin{pmatrix}
 0&-1&-1&4\\0&0&-1&2\\1&1&1&-6\\0&0&0&1
 \end{pmatrix}.
\end{equation}
It is obtained from the previously verified determinant-one
permutation basis by adding $-2(u_0+u_1+u_2)$ to the last
column, so its determinant is one. Since $T$ fixes both
$e_4$ and $U$, it fixes $v_4$ and permutes the other cells.
These are therefore genuine transported product cells, obtained
by an integral diffeomorphism, in every degree.

In these cells the maps in \eqref{marked-local:eq:marked-native-FK} are
$[y,3t]$ and $[y,3t+v]$. Their interval subdivisions are exactly
those in Theorem~\ref{marked-local:thm:positive-chains}, which proves all
formulas in \eqref{marked-local:eq:marked-final-chains} and all their chain
identities. Proposition~\ref{marked-local:prop:positive-splitting} gives the
section and kernel of the full boundary map. In particular the
degree-four map and all adjacent prism terms are included.
\end{proof}

The original twist has not been changed: in this cell basis
$e_4/3=(v_4+2U)/3$. Its sum with the original central translation
$U/3$ is $v_4/3+U$, which equals $v_4/3$ as a torus translation.
Both contributions are needed for the integral comparison.

For a homology class written in the original period basis, the
same maps are explicit as well. In degree $n$, change its exterior
coordinates by $\bigwedge^nV^{-1}$, express the $u$-pairs in
$(w_0,w_1,w_2)$, and move $v_4$ to the front in the suspended
terms before applying \eqref{marked-local:eq:marked-final-chains}. Moving it
past $n-1$ fibre vectors contributes $(-1)^{n-1}$.
Thus this is an integral change of actual geometric
cells, with no rational replacement of the lattice.
For example, write $c=[u_0]=[u_1]=[u_2]$ and let $h$ be the
positive core interval in $H_1(G_P;\mathbb Z)$. Then
\begin{equation}\label{marked-local:eq:marked-degree-one}
 f_{0*}(e_1)=f_{0*}(e_2)=0,\qquad
 f_{0*}(e_3)=c,\qquad f_{0*}(e_4)=6c+3h,
 \qquad k_{0*}(\gamma_0)=h.
\end{equation}
Indeed $e_1=u_0-u_1$, $e_2=u_1-u_2$, $e_3=u_0$,
and $e_4=v_4+2U$. This gives a deciding distinction from the
zero-origin descent in the unchanged original marking.
At the level of the fundamental group, three traversals of
the marked meridian satisfy
\begin{equation}\label{marked-local:eq:marked-meridian-relation}
                   \gamma_0^3=f_{0\#}(e_4-2\delta).
\end{equation}
To verify it directly, lift \eqref{marked-local:eq:marked-native-j} to the
universal vector coordinates. Its endpoint is the affine deck
generator followed by translation by $-\delta$. Cubing adds
$e_4+\delta-3\delta=e_4-2\delta$. Equivalently in $G_P$, the
marked loop is its positive core interval, while the $v_4$ fibre
loop covers three such intervals. The image of $\delta$ is
$3c$ and is not zero in this local filling. Consequently the
local relation cannot be replaced by $\gamma_0^3=f_{0\#}(e_4)$.

For the higher degrees, abbreviate $e_i\wedge e_j$ by $e_{ij}$
and similarly for three or four indices. In the ordered domain
bases $(e_{12},e_{13},e_{14},e_{23},e_{24},e_{34})$ and
$(e_{123},e_{124},e_{134},e_{234})$, the degree-two and
degree-three maps are, respectively,
\begin{equation}\label{marked-local:eq:marked-higher-matrices}
 f_{0*,2}=\begin{pmatrix}3&1&0&-2&0&0\\0&0&0&0&0&-1\end{pmatrix},
 \qquad
 f_{0*,3}=\begin{pmatrix}1&6&2&-4\\0&3&1&-2\end{pmatrix}.
\end{equation}
The target bases are $([w_0],\sigma U)$ and
$(\omega,\sigma W)$ in the core complex. For example,
$e_{12}=w_0+w_1+w_2$ gives the first coefficient $3$.
The term $e_3\wedge v_4=-v_4\wedge u_0$ gives the last
coefficient $-1$ in degree two. In degree three,
$e_{124}=v_4\wedge e_{12}+2e_{12}\wedge U$
gives $3\sigma W+6\omega$. Expanding the other columns
in the same stated basis gives \eqref{marked-local:eq:marked-higher-matrices}.
Both maps are integrally surjective: $e_{13},-e_{34}$ map to
the two target generators in degree two, and
$e_{123},e_{134}-2e_{123}$ do so in degree three.
Their kernels have the following full integral bases:
\[
 \begin{array}{c|l}
 2&e_{12}-3e_{13},\quad e_{14},\quad 2e_{13}+e_{23},\quad e_{24},\\
 3&e_{124}-3e_{134},\quad 2e_{134}+e_{234}.
 \end{array}
\]
Indeed the two equations in each displayed matrix solve for the
remaining coordinates with integer coefficients. In degree four,
$e_{1234}=-v_4\wedge\omega$, so its image is $-3\sigma\omega$.
The covering has geometric degree $+3$: the interval-first target
cell $\sigma\omega$ is the negative of the fibre-first core
orientation. These formulas retain both the signs and the complete
integral images in the original periods.

All homology groups, kernel ranks, and cokernels in
Theorem~\ref{marked-local:thm:homology} remain valid in the stated cell basis.
The invariant degree-one core covector pulls back to
$e_3^*+6e_4^*$, and the core-circle covector pulls back to
$3e_4^*$. Their image is still
$\mathbb Ze_3^*\oplus3\mathbb Ze_4^*$. The higher-degree
specialization indices remain $1,1,3$ in degrees two, three,
and four, by the explicit unimodular change of cells.
Equal indices therefore do not imply equal marked attachment maps.

This local conclusion depends on the smooth line-bundle reduction,
its commuting lifted translation, and its stated period marking.
Within those hypotheses the conormal comparison constructs the
affine deck action and the positive boundary map. Comparisons with
a common chain complex carrying the other local monodromies still
require explicit maps and homotopies. The other finite-order filling,
the cusp attachments, and their global integral gluing have not been
constructed here. No conclusion about the homology of their union
or a complex structure on a sphere follows from this local calculation.

